% ================================================================
%  sesion15.tex  —  Sesión 15 de 20
%  Comparar gráficas
%  Bloque 3: Datos y gráficas · Matemáticas 1.º ESO
% ================================================================
\documentclass[aspectratio=169, 12pt, spanish]{beamer}
\usepackage{almeria-beamer}

\renewcommand{\SessionNumber}{15}
\renewcommand{\SessionTitle}{Comparar gráficas}
\renewcommand{\SessionName}{Sesión 15}
\renewcommand{\SessionObjective}{%
  Superponer y comparar dos gráficas sobre el mismo sistema de ejes,
  identificar correlaciones positivas y negativas, y extraer
  conclusiones comparativas sobre datos urbanos de Almería.}
\renewcommand{\BlockName}{Bloque 3: Datos y gráficas}

\begin{document}

\begin{frame}[plain]\titlepage\end{frame}

\begin{frame}{Índice}
  \begin{columns}[t]
    \begin{column}{0.5\textwidth}
      \begin{enumerate}
        \item ¿Por qué comparar gráficas?
        \item Tipos de correlación
        \item Temperatura vs.\ turismo en Almería
        \item Lluvia vs.\ aforo en playas
        \item Correlación $\neq$ causalidad
      \end{enumerate}
    \end{column}
    \begin{column}{0.5\textwidth}
      \begin{enumerate}\setcounter{enumi}{5}
        \item Cómo «miente» una gráfica
        \item Ejercicios multinivel (Bloom)
        \item Evidencia del día
        \item Error frecuente
        \item Resumen
      \end{enumerate}
    \end{column}
  \end{columns}
\end{frame}

% ---------------------------------------------------------------
\seccionframe{1. ¿Por qué comparar gráficas?}
% ---------------------------------------------------------------

\begin{frame}{La comparación revela relaciones}
  \begin{columns}[c]
    \begin{column}{0.52\textwidth}
      \begin{teoriabox}
        Superponer dos gráficas en los mismos ejes permite:
        \begin{itemize}
          \item Detectar si dos variables
                \textbf{suben y bajan juntas}
                (correlación positiva).
          \item Detectar si una \textbf{sube cuando
                la otra baja} (correlación negativa).
          \item Identificar \textbf{retrasos} (lag):
                una variable responde a la otra con demora.
          \item Descubrir períodos de
                \textbf{comportamiento anómalo}.
        \end{itemize}
      \end{teoriabox}
    \end{column}
    \begin{column}{0.44\textwidth}
      \begin{tikzpicture}[every node/.style={font=\scriptsize}]
        \begin{axis}[
          almeria line,
          width=6.5cm, height=4.5cm,
          xmin=1, xmax=6, ymin=0, ymax=11,
          xtick={1,...,6},
          xticklabels={E,F,M,A,M,J},
          ylabel={Valor (escalado)},
          xlabel={Mes},
          legend pos=north west,
          legend style={font=\tiny},
        ]
          \addplot[AlmTerra, line width=2pt, mark=*]
            coordinates {(1,2)(2,3)(3,5)(4,7)(5,9)(6,10)};
          \addlegendentry{Variable A}
          \addplot[AlmAzulM, line width=2pt, mark=square*]
            coordinates {(1,8)(2,7)(3,6)(4,4)(5,3)(6,2)};
          \addlegendentry{Variable B}
        \end{axis}
      \end{tikzpicture}
      \begin{center}
        \small\color{AlmGris}
        A sube, B baja → correlación \textbf{negativa}
      \end{center}
    \end{column}
  \end{columns}
\end{frame}

% ---------------------------------------------------------------
\seccionframe{2. Tipos de correlación}
% ---------------------------------------------------------------

\begin{frame}{Correlación: positiva, negativa o nula}
  \begin{columns}[T]
    \begin{column}{0.31\textwidth}
    \centering
    {\color{AlmVerde}\bfseries Correlación positiva}\\[4pt]
    \begin{tikzpicture}
      \begin{axis}[
        width=4.5cm, height=3.8cm,
        axis lines=left, xmin=0, xmax=6, ymin=0, ymax=10,
        xtick=\empty, ytick=\empty,
        tick style={draw=none},
      ]
        \addplot[AlmVerde, line width=2pt, no marks]
          coordinates {(0,1)(5,9)};
        \addplot[only marks, mark=*, AlmVerde!80, mark size=2pt]
          coordinates {(1,1.8)(2,3.5)(3,5.1)(4,6.8)(5,8.5)};
      \end{axis}
    \end{tikzpicture}\\[4pt]
    \footnotesize Cuando $x$ $\uparrow$, también $y$ $\uparrow$\\
    Ejemplo: temperatura $\uparrow$ → turistas $\uparrow$
    \end{column}
    \begin{column}{0.31\textwidth}
    \centering
    {\color{AlmTerra}\bfseries Correlación negativa}\\[4pt]
    \begin{tikzpicture}
      \begin{axis}[
        width=4.5cm, height=3.8cm,
        axis lines=left, xmin=0, xmax=6, ymin=0, ymax=10,
        xtick=\empty, ytick=\empty,
        tick style={draw=none},
      ]
        \addplot[AlmTerra, line width=2pt, no marks]
          coordinates {(0,9)(5,1)};
        \addplot[only marks, mark=*, AlmTerra!80, mark size=2pt]
          coordinates {(1,8.2)(2,6.5)(3,4.9)(4,3.2)(5,1.5)};
      \end{axis}
    \end{tikzpicture}\\[4pt]
    \footnotesize Cuando $x$ $\uparrow$, $y$ $\downarrow$\\
    Ejemplo: lluvia $\uparrow$ → aforo playa $\downarrow$
    \end{column}
    \begin{column}{0.31\textwidth}
    \centering
    {\color{AlmGris}\bfseries Sin correlación}\\[4pt]
    \begin{tikzpicture}
      \begin{axis}[
        width=4.5cm, height=3.8cm,
        axis lines=left, xmin=0, xmax=6, ymin=0, ymax=10,
        xtick=\empty, ytick=\empty,
        tick style={draw=none},
      ]
        \addplot[only marks, mark=*, AlmGris, mark size=2pt]
          coordinates {
            (0.5,2)(1,8)(1.5,4)(2,6)(2.5,1)(3,9)
            (3.5,3)(4,7)(4.5,2)(5,5)(5.5,8)
          };
      \end{axis}
    \end{tikzpicture}\\[4pt]
    \footnotesize Los puntos no siguen ningún patrón\\
    Ejemplo: n.º de farolas $\leftrightarrow$ turistas
    \end{column}
  \end{columns}
\end{frame}

% ---------------------------------------------------------------
\seccionframe{3. Temperatura vs.\ turismo}
% ---------------------------------------------------------------

\begin{frame}{Temperatura y turismo en Almería — dos ejes}
  \begin{columns}[c]
    \begin{column}{0.62\textwidth}
      \begin{tikzpicture}
        \begin{axis}[
          width=10cm, height=5.8cm,
          axis y line*=left,
          axis x line=bottom,
          xmin=1, xmax=12, ymin=8, ymax=33,
          xtick={1,...,12},
          xticklabels={E,F,M,A,M,J,J,A,S,O,N,D},
          ylabel={\color{AlmTerra}Temperatura (°C)},
          xlabel={Mes},
          tick label style={font=\tiny},
          label style={font=\small},
          title={Temperatura (°C) y turistas ($\times 10^4$)},
          title style={font=\small\bfseries, color=AlmAzul},
          grid=both,
          grid style={line width=0.3pt, AlmAzulC!25},
        ]
          \addplot[AlmTerra, line width=2.5pt, mark=*,
                   mark options={fill=AlmTerra}]
            coordinates {
              (1,12)(2,13)(3,16)(4,19)(5,23)(6,27)
              (7,30)(8,30)(9,26)(10,21)(11,16)(12,13)
            };
          \addlegendentry{\footnotesize Temperatura}
        \end{axis}
        \begin{axis}[
          width=10cm, height=5.8cm,
          axis y line*=right,
          axis x line=none,
          xmin=1, xmax=12, ymin=0, ymax=80,
          ylabel={\color{AlmAzulM}Turistas ($\times 10^4$/mes)},
          tick label style={font=\tiny},
          label style={font=\small},
          ytick={0,20,40,60,80},
        ]
          \addplot[AlmAzulM, line width=2.5pt, mark=square*,
                   mark options={fill=AlmAzulM}, dashed]
            coordinates {
              (1,6)(2,6.5)(3,8)(4,12)(5,15)(6,25)
              (7,50)(8,60)(9,30)(10,20)(11,10)(12,8)
            };
          \addlegendentry{\footnotesize Turistas}
        \end{axis}
      \end{tikzpicture}
    \end{column}
    \begin{column}{0.34\textwidth}
      \textbf{Observaciones:}
      \begin{itemize}
        \item Ambas curvas tienen el \textbf{pico en julio--agosto}.
        \item Ambas tienen el \textbf{mínimo en enero}.
        \item Cuando la temperatura $\uparrow$,
              los turistas también $\uparrow$.
        \item \textbf{Correlación positiva} clara.
      \end{itemize}
      \vspace{6pt}
      \begin{notabox}[Doble eje $y$]
        \footnotesize
        Se usan dos ejes $y$ cuando las
        variables tienen \textbf{escalas muy
        diferentes} (°C vs.\ decenas de miles).
      \end{notabox}
    \end{column}
  \end{columns}
\end{frame}

% ---------------------------------------------------------------
\seccionframe{4. Lluvia vs.\ aforo en playas}
% ---------------------------------------------------------------

\begin{frame}{Lluvia y aforo de playas — correlación negativa}
  \begin{columns}[c]
    \begin{column}{0.58\textwidth}
      \begin{tikzpicture}
        \begin{axis}[
          almeria bar,
          width=9.5cm, height=5.2cm,
          symbolic x coords={E,F,M,A,M,J,J,A,S,O,N,D},
          xtick=data,
          ylabel={Lluvia (mm) / Aforo (\%)},
          xlabel={Mes},
          title={Lluvia mensual (barras) y aforo de playa (línea)},
          ymin=0, ymax=100,
          bar width=0.38cm,
          legend pos=north west,
          legend style={font=\tiny},
          nodes near coords=\empty,
        ]
          \addplot[fill=AlmAzulC!60, fill opacity=0.8]
            coordinates {
              (E,25)(F,22)(M,18)(A,12)(M,8)(J,3)
              (J,1)(A,1)(S,8)(O,20)(N,28)(D,30)
            };
          \addlegendentry{Lluvia (mm)}
          \addplot[AlmTerra, line width=2.5pt, mark=*,
                   mark options={fill=AlmTerra}]
            coordinates {
              (E,15)(F,18)(M,22)(A,35)(M,50)(J,70)
              (J,95)(A,98)(S,65)(O,35)(N,18)(D,12)
            };
          \addlegendentry{Aforo playa (\%)}
        \end{axis}
      \end{tikzpicture}
    \end{column}
    \begin{column}{0.38\textwidth}
      \textbf{Análisis:}
      \begin{itemize}
        \item Lluvia alta (otoño--invierno) →
              aforo de playa \textbf{bajo}.
        \item Lluvia baja (verano) →
              aforo de playa \textbf{alto}.
        \item \textbf{Correlación negativa}: cuando
              una variable sube, la otra baja.
      \end{itemize}
      \vspace{6pt}
      \begin{analizarbox}
        \footnotesize
        ¿La lluvia es la \textit{causa} de que
        las playas estén vacías, o existen
        otras variables explicativas?
        (temperatura, vacaciones, etc.)
      \end{analizarbox}
    \end{column}
  \end{columns}
\end{frame}

% ---------------------------------------------------------------
\seccionframe{5. Correlación $\neq$ Causalidad}
% ---------------------------------------------------------------

\begin{frame}{¡Atención! Correlación no implica causalidad}
  \begin{columns}[T]
    \begin{column}{0.55\textwidth}
      \begin{errorbox}
        \textbf{Error de razonamiento frecuente:}\\[4pt]
        «Como la temperatura y el turismo suben juntos,
        \textit{la temperatura es la causa del turismo.}»\\[6pt]
        Pero también hay \textbf{otra variable oculta}:
        las \textit{vacaciones escolares} coinciden con
        el verano. Son las vacaciones (no solo el calor)
        las que permiten a las familias viajar.
      \end{errorbox}
      \vspace{6pt}
      \begin{notabox}[Variable oculta o confusora]
        \footnotesize
        Una \textbf{variable confusora} es aquella que
        influye en $x$ e $y$ simultáneamente, haciendo
        que parezca que $x$ causa $y$ cuando en realidad
        ambas están causadas por una tercera variable.
      \end{notabox}
    \end{column}
    \begin{column}{0.42\textwidth}
      \textbf{Ejemplo clásico:}\\[6pt]
      \begin{center}
      \begin{tikzpicture}[
        nodo/.style={
          rounded corners=4pt, draw=AlmAzulM,
          fill=AlmFondo, font=\small, inner sep=5pt,
          text width=2.5cm, align=center
        },
        causa/.style={->, AlmTerra, line width=1.5pt}
      ]
        \node[nodo] (V) at (0,2.5) {Vacaciones\\de verano};
        \node[nodo] (T) at (-2,0) {Temperatura\\alta (°C)};
        \node[nodo] (Tu) at (2,0) {Turistas\\$\uparrow$};
        \draw[causa] (V) -- (T);
        \draw[causa] (V) -- (Tu);
        \draw[AlmGris, dashed, ->, line width=1pt] (T) -- (Tu)
          node[midway, below, font=\tiny, color=AlmGris]
          {correlación\,$\neq$\,causa};
      \end{tikzpicture}
      \end{center}
      \vspace{4pt}
      \footnotesize\color{AlmGris}
      Las vacaciones causan TANTO que haya mucha gente
      en la playa COMO que observemos días calurosos
      en los registros turísticos.
    \end{column}
  \end{columns}
\end{frame}

% ---------------------------------------------------------------
\seccionframe{6. Cómo «miente» una gráfica}
% ---------------------------------------------------------------

\begin{frame}{La escala puede engañar}
  \begin{columns}[T]
    \begin{column}{0.48\textwidth}
      \centering
      {\color{AlmTerra}\bfseries\small Gráfica «alarmista»}\\
      (eje $y$ truncado: 0 omitido)\\[4pt]
      \begin{tikzpicture}
        \begin{axis}[
          width=6cm, height=4.2cm,
          axis lines=left,
          xmin=2019.5, xmax=2022.5,
          ymin=2.7, ymax=3.1,
          xtick={2020,2021,2022},
          ytick={2.7,2.8,2.9,3.0,3.1},
          xlabel={Año}, ylabel={Mill.},
          tick label style={font=\tiny},
          label style={font=\tiny},
          title style={font=\tiny},
          title={¡Caída catastrófica!},
        ]
          \addplot[AlmTerra, line width=2.5pt, mark=*]
            coordinates {(2020,2.75)(2021,2.85)(2022,3.0)};
        \end{axis}
      \end{tikzpicture}
    \end{column}
    \begin{column}{0.48\textwidth}
      \centering
      {\color{AlmVerde}\bfseries\small Gráfica honesta}\\
      (eje $y$ desde 0)\\[4pt]
      \begin{tikzpicture}
        \begin{axis}[
          width=6cm, height=4.2cm,
          axis lines=left,
          xmin=2019.5, xmax=2022.5,
          ymin=0, ymax=3.5,
          xtick={2020,2021,2022},
          ytick={0,1,2,3},
          xlabel={Año}, ylabel={Mill.},
          tick label style={font=\tiny},
          label style={font=\tiny},
          title style={font=\tiny},
          title={Crecimiento moderado},
        ]
          \addplot[AlmVerde, line width=2.5pt, mark=*]
            coordinates {(2020,2.75)(2021,2.85)(2022,3.0)};
        \end{axis}
      \end{tikzpicture}
    \end{column}
  \end{columns}
  \vspace{6pt}
  \begin{center}
    \small ¡Los mismos datos, escalas diferentes,
    impresiones \textbf{completamente distintas}!
    Siempre comprueba si el eje $y$ empieza en $0$.
  \end{center}
\end{frame}

% ---------------------------------------------------------------
\seccionframe{7. Ejercicios multinivel}
% ---------------------------------------------------------------

\begin{frame}{Ejercicios multinivel}
  \begin{columns}[T]
    \begin{column}{0.31\textwidth}
      \begin{analizarbox}
        \footnotesize\dificultad{2}\\[4pt]
        Tienes un gráfico dual: temperatura (°C)
        y ventas de helados en Almería por mes.\\[4pt]
        \textbf{Analiza:}\\
        (a) Tipo de correlación.\\
        (b) ¿Cuándo coinciden los picos?\\
        (c) ¿Hay un desfase (lag) entre ambas?\\
        Argumenta en 5 líneas.
      \end{analizarbox}
    \end{column}
    \begin{column}{0.31\textwidth}
      \begin{evaluarbox}
        \footnotesize\dificultad{3}\\[4pt]
        Tienes tres comparaciones:\\[3pt]
        A: temperatura vs.\ turistas\\
        B: vacaciones escolares vs.\ turistas\\
        C: precio vuelos vs.\ turistas\\[4pt]
        \textbf{Evalúa} cuál de los tres factores
        explica mejor el turismo en Almería.
        ¿Son independientes o están relacionados
        entre sí? Escribe un argumento de 8 líneas.
      \end{evaluarbox}
    \end{column}
    \begin{column}{0.31\textwidth}
      \begin{crearbox}
        \footnotesize\dificultad{4}\\[4pt]
        Inventa \textbf{dos variables urbanas}
        de Almería que tengan:
        \begin{itemize}
          \item Una correlación \textbf{positiva}
          \item Una correlación \textbf{negativa}
        \end{itemize}
        Dibuja a mano un gráfico esquemático
        de cada par. Explica la posible variable
        confusora en cada caso.
      \end{crearbox}
    \end{column}
  \end{columns}
\end{frame}

% ---------------------------------------------------------------
\begin{frame}{Evidencia del día}
  \begin{evidenciabox}
    Escribe un \textbf{análisis comparativo} de dos gráficas de
    datos de Almería. Debe incluir:
    \begin{enumerate}
      \item Identificación del tipo de correlación (positiva/negativa/nula)
            con justificación basada en los valores leídos
      \item Al menos un punto donde ambas variables cambian juntas
            y uno donde se separan — explica por qué
      \item Discusión de si la correlación implica causalidad:
            ¿existe una variable oculta?
      \item Conclusión de 4--5 líneas sobre la relación entre
            las dos variables en el contexto de Almería
    \end{enumerate}
  \end{evidenciabox}
\end{frame}

% ---------------------------------------------------------------
\begin{frame}{¡Cuidado con este error!}
  \begin{errorbox}
    \textbf{Correlación $\neq$ Causalidad}\\[6pt]
    \incorrecto\ «El turismo aumenta porque hace más calor.
    Por tanto, si calentamos la ciudad artificialmente,
    vendrán más turistas.»\\[6pt]
    \correcto\ La correlación entre temperatura y turismo se debe en
    parte a las vacaciones de verano, al sol en la playa, y a
    patrones culturales. El calor \textit{solo} no causaría más turismo.
    Identificar la variable confusora es fundamental para
    razonar correctamente.
  \end{errorbox}
\end{frame}

% ---------------------------------------------------------------
\begin{frame}{Resumen de la sesión 15}
  \begin{resumenbox}
    \begin{itemize}
      \item Comparar gráficas permite detectar
            \textbf{correlaciones} entre variables.
      \item \textbf{Correlación positiva}: ambas variables
            suben/bajan juntas.
      \item \textbf{Correlación negativa}: una sube cuando
            la otra baja.
      \item Se usan \textbf{dos ejes $y$} cuando las escalas
            son muy distintas.
      \item \textbf{Correlación $\neq$ causalidad}:
            siempre busca posibles variables confusoras.
      \item Una gráfica puede \textbf{engañar} si el eje $y$
            no empieza en $0$.
    \end{itemize}
  \end{resumenbox}
  \vspace{6pt}
  \begin{center}
    {\small\color{AlmAzulM}\faArrowCircleRight\enskip
    \textbf{Sesión 16:} Relaciones funcionales y proporcionalidad directa}
  \end{center}
\end{frame}

\end{document}
